
The point-to-point synchronization routines block or test until a symmetric data
object on the local \ac{PE} satisfies a specified condition.

Each of these routines is also provided in thread group (\texttt{\_tg}) and
hardware thread group (\texttt{\_htg}) variants, as defined in
\cref{subsec:gc_comm_ops}. In these collective variants, all participating
threads return together once the condition is satisfied, and any return value is
the same across the group.

The \VAR{ivar} or \VAR{sig\_addr} data object shall reside in the \ac{GPU}
symmetric heap. Using local \ac{GPU} memory, local CPU memory, or memory from the
default symmetric heap results in \textbf{undefined behavior}.

\subsubsubsection{\textbf{Wait-Until}}\label{subsec:gc_wait_until}

Blocks until a symmetric data object on the local \ac{PE} satisfies a condition.

\begin{apidefinition}
\begin{Csynopsis}
SHMEMG_DEVICE void shmemg_TYPENAME_wait_until(TYPE *ivar, int cmp,
    TYPE cmp_value);
\end{Csynopsis}

Thread group variant:
\begin{CsynopsisCol}
SHMEMG_DEVICE void shmemg_TYPENAME_wait_until_tg(TYPE *ivar, int cmp,
    TYPE cmp_value);
\end{CsynopsisCol}

Hardware thread group variant:
\begin{CsynopsisCol}
SHMEMG_DEVICE void shmemg_TYPENAME_wait_until_htg(TYPE *ivar, int cmp,
    TYPE cmp_value);
\end{CsynopsisCol}

where \TYPE{} is one of the point-to-point synchronization types defined in
\openshmem 1.6, and \TYPENAME{} is the corresponding type-name component of the
routine name.

\begin{apiarguments}
\apiargument{IN}{ivar}{Symmetric address of the source data object. Shall reside
in the \ac{GPU} symmetric heap.}
\apiargument{IN}{cmp}{The comparison operator.}
\apiargument{IN}{cmp\_value}{The value against which \VAR{ivar} will be compared.}
\end{apiarguments}

\apidescription{
Callable from within a \ac{GPU} kernel.
\FUNC{shmemg\_wait\_until} blocks until the value of \VAR{ivar} at the calling
\ac{PE} satisfies the condition implied by \VAR{cmp} and \VAR{cmp\_value}. The
supported values of \VAR{cmp} are: \CONST{SHMEM\_CMP\_EQ},
\CONST{SHMEM\_CMP\_NE}, \CONST{SHMEM\_CMP\_GT}, \CONST{SHMEM\_CMP\_GE},
\CONST{SHMEM\_CMP\_LT}, and \CONST{SHMEM\_CMP\_LE}.

The routine shall not return until the condition is satisfied. If no \ac{PE} ever
updates \VAR{ivar} to satisfy the condition, the routine shall not return.
}

\apireturnvalues{
None.
}
\end{apidefinition}

\subsubsubsection{\textbf{Test}}\label{subsec:gc_test}

Tests whether a symmetric data object on the local \ac{PE} satisfies a
condition.

\begin{apidefinition}
\begin{Csynopsis}
SHMEMG_DEVICE int shmemg_TYPENAME_test(TYPE *ivar, int cmp, TYPE cmp_value);
\end{Csynopsis}

Thread group variant:
\begin{CsynopsisCol}
SHMEMG_DEVICE int shmemg_TYPENAME_test_tg(TYPE *ivar, int cmp, TYPE cmp_value);
\end{CsynopsisCol}

Hardware thread group variant:
\begin{CsynopsisCol}
SHMEMG_DEVICE int shmemg_TYPENAME_test_htg(TYPE *ivar, int cmp, TYPE cmp_value);
\end{CsynopsisCol}

where \TYPE{} is one of the point-to-point synchronization types defined in
\openshmem 1.6, and \TYPENAME{} is the corresponding type-name component of the
routine name.

\begin{apiarguments}
\apiargument{IN}{ivar}{Symmetric address of the source data object. Shall reside
in the \ac{GPU} symmetric heap.}
\apiargument{IN}{cmp}{The comparison operator.}
\apiargument{IN}{cmp\_value}{The value against which \VAR{ivar} will be compared.}
\end{apiarguments}

\apidescription{
Callable from within a \ac{GPU} kernel.
\FUNC{shmemg\_test} tests whether the value of \VAR{ivar} at the calling \ac{PE}
satisfies the condition implied by \VAR{cmp} and \VAR{cmp\_value}. The test is
performed once and returns without blocking.
}

\apireturnvalues{
1 if the condition is satisfied; 0 otherwise.
}
\end{apidefinition}

\subsubsubsection{\textbf{Signal-Wait-Until}}\label{subsec:gc_signal_wait_until}

Blocks until a local signal data object satisfies a condition.

\begin{apidefinition}
\begin{Csynopsis}
SHMEMG_DEVICE uint64_t shmemg_signal_wait_until(uint64_t *sig_addr,
    int cmp, uint64_t cmp_value);
\end{Csynopsis}

Thread group variant:
\begin{CsynopsisCol}
SHMEMG_DEVICE uint64_t shmemg_signal_wait_until_tg(uint64_t *sig_addr,
    int cmp, uint64_t cmp_value);
\end{CsynopsisCol}

Hardware thread group variant:
\begin{CsynopsisCol}
SHMEMG_DEVICE uint64_t shmemg_signal_wait_until_htg(uint64_t *sig_addr,
    int cmp, uint64_t cmp_value);
\end{CsynopsisCol}

\begin{apiarguments}
\apiargument{IN}{sig\_addr}{Local address of the remotely accessible signal
variable. Shall reside in the \ac{GPU} symmetric heap.}
\apiargument{IN}{cmp}{The comparison operator.}
\apiargument{IN}{cmp\_value}{The value against which \VAR{sig\_addr} will be
compared.}
\end{apiarguments}

\apidescription{
Callable from within a \ac{GPU} kernel.
\FUNC{shmemg\_signal\_wait\_until} blocks until the value of the signal data
object at \VAR{sig\_addr} satisfies the condition implied by \VAR{cmp} and
\VAR{cmp\_value}. The routine shall not return until the condition is satisfied.

Access to \VAR{sig\_addr} satisfies the atomicity guarantees described in the
signaling operations introduction.
}

\apireturnvalues{
The value of the signal data object at \VAR{sig\_addr} that satisfies the wait
condition.
}
\end{apidefinition}

\paragraph*{\textbf{Example}}\mbox{}\par\nobreak
The following program performs a point-to-point handshake. \ac{PE}~0 sets a flag
on \ac{PE}~1 with a put, and \ac{PE}~1 blocks inside its kernel with
\FUNC{shmemg\_long\_wait\_until} until the flag is observed. A nonblocking
\FUNC{shmemg\_long\_test} could be used instead to poll without blocking.

\lstinputlisting[style=SourceListingsDefault]{example_code/gc_p2p_example.c}

